\chapter{Інженерія виявлення: правила, індикатори та безфайлові загрози}
\chaptershort{Інженерія виявлення}
\label{chap:ch11}

Методи виявлення (розділи~\ref{chap:ch06}--\ref{chap:ch10}) стають дієвими лише тоді, коли втілені в конкретні, супроводжувані, вимірювані \term{правила виявлення}. \term{Інженерія виявлення} (detection engineering) --- це систематична дисципліна створення, тестування, розгортання й супроводу логіки виявлення, що розглядає правила як код: із контролем версій, тестуванням, метриками якості й життєвим циклом. Цей розділ поєднує попередні методи в практику: він розглядає мови й рівні правил, «піраміду болю» та ієрархію індикаторів, виявлення безфайлових атак і зловживання легітимними інструментами (living-off-the-land), життєвий цикл і метрики виявлення, а також покриття технік ATT\&CK, і завершується кейсом розроблення виявлення для конкретної техніки. Результати живлять захист кінцевих точок (розділ~\ref{chap:ch10}), розвідку загроз (розділ~\ref{chap:ch12}), полювання за загрозами (розділ~\ref{chap:ch13}) і реагування (розділ~\ref{chap:ch16}).

\section{Правила виявлення як код}\label{sec:ch11-rules}

Інженерія виявлення запозичує практики розробки ПЗ: правило має версію, автора, опис, тести, метрики й життєвий цикл. Це протиставляється стихійному накопиченню несупроводжуваних правил, що з часом деградують (хибні тривоги, застарілість). Правила класифікують за рівнем даних, на якому вони працюють: \emph{файлові} (YARA, розділ~\ref{chap:ch06}) --- над вмістом і структурою файлів; \emph{подієві/хостові} (Sigma~\cite{sigma}) --- над журналами подій кінцевих точок (розділ~\ref{chap:ch10}); \emph{мережеві} (Snort~\cite{ch11:snort}, Suricata~\cite{ch11:suricata}) --- над мережевим трафіком. Кожен рівень бачить свій зріз активності, тому повне покриття вимагає правил на всіх рівнях.

\begin{tikzfig}
\begin{tikzpicture}[font=\scriptsize,
  b/.style={draw,rounded corners=2pt,minimum height=6mm,text width=22mm,align=center,font=\tiny},
  >={Stealth[length=1.5mm]}]
\node[b] (f) at (0,0) {файлові (YARA)};
\node[b] (h) at (0,-1.1) {хостові (Sigma)};
\node[b] (n) at (0,-2.2) {мережеві (Suricata)};
\node[b] (corr) at (3.4,-1.1) {кореляція / SIEM};
\node[b] (alert) at (6.6,-1.1) {тривога + контекст};
\draw[->] (f)--(corr); \draw[->] (h)--(corr); \draw[->] (n)--(corr);
\draw[->] (corr)--(alert);
\end{tikzpicture}
\figcap{Рис.~\thechapter.1 --- Рівні правил виявлення й їх злиття в кореляційній системі}
\end{tikzfig}

Якість правила визначають метриками детектора (розділ~\ref{chap:ch03}): повнотою, точністю, частотою хибних тривог. Інженерія виявлення додає й операційні характеристики: вартість обчислення (чи встигає правило в реальному часі), супроводжуваність (чи зрозуміле воно, чи легко оновити), стійкість до ухилення (наскільки легко обійти). Зріле правило документує не лише \emph{що} воно шукає, а й \emph{чому}, на яку техніку спрямоване та які відомі способи обходу.

Перенесення практик розробки ПЗ на виявлення (detection-as-code) має конкретні наслідки. \emph{Контроль версій} зберігає історію змін правила, дозволяє відкотити невдале оновлення й пов'язати зміну з причиною (нова загроза, хибні тривоги). \emph{Автоматизоване тестування} прогонить правило проти набору злоякісних зразків (перевірка повноти) і проти «чистого» еталонного трафіку (перевірка на хибні тривоги) перед кожним розгортанням --- так само, як модульні тести захищають код від регресій. \emph{Рецензування} (code review) правил іншими інженерами виявляє логічні помилки, надмірну специфічність і прогалини. \emph{Безперервна інтеграція} автоматизує цей конвеєр. Без цих практик бібліотека правил великої організації швидко стає некерованою: ніхто не пам'ятає, навіщо існує те чи інше правило, які з них дублюються, а які давно не спрацьовують, --- і сукупна якість виявлення непомітно деградує.

Важливо також розрізняти \emph{атомарні} та \emph{кореляційні} правила. Атомарне правило спрацьовує на одній події чи ознаці (підозрілий командний рядок, звернення до відомого шкідливого домену). Кореляційне правило поєднує кілька подій у часі та за сутностями (той самий вузол, той самий користувач) --- наприклад, «фішинговий лист, а згодом на тому ж вузлі аномальний процес і вихідне з'єднання». Атомарні правила прості, швидкі, але шумні; кореляційні --- точніші й стійкіші, але складніші в розробленні й потребують якісної, збагаченої телеметрії (розділ~\ref{chap:ch10}). Зріла програма виявлення поєднує обидва типи.

\section{Піраміда болю та ієрархія індикаторів}\label{sec:ch11-pyramid}

Ключова концепція інженерії виявлення --- \term{піраміда болю}~\cite{pyramid}, що впорядковує індикатори за тим, наскільки \emph{болісно} зловмиснику їх змінити (а отже, наскільки стійке виявлення на їх основі). Знизу вгору: геші файлів (змінити тривіально --- один байт), IP-адреси (легко --- змінити інфраструктуру), доменні імена (трохи складніше), мережеві й хостові артефакти (складніше), інструменти (ще складніше --- потрібно переписати чи замінити), і на вершині --- \term{тактики, техніки й процедури} (TTP), змінити які найважче, бо це вимагає змінити сам спосіб дій. Цю ієрархію узагальнено в табл.~\thechapter.1.

\tabhead{Таблиця~\thechapter.1 --- Піраміда болю: ієрархія індикаторів за стійкістю}

\begin{center}\footnotesize
\begin{tabularx}{\linewidth}{L{22mm} Y L{18mm}}
\toprule
Рівень & Індикатор & Вартість зміни \\
\midrule
Нижній & геші файлів & тривіальна \\
 & IP-адреси, домени & низька \\
Середній & артефакти хоста/мережі & помірна \\
 & інструменти & висока \\
Вершина & TTP (техніки, процедури) & дуже висока \\
\bottomrule
\end{tabularx}
\end{center}

Практичний висновок: виявлення на нижніх рівнях (геші, адреси) дешеве й точне, але нестійке --- зловмисник обходить його миттєво. Виявлення на вершині (поведінкові патерни TTP) дороге в розробленні, але стійке, бо змушує зловмисника змінити спосіб дій. Тому зріла стратегія виявлення інвестує у \emph{поведінкові} правила над техніками, тоді як індикатори нижніх рівнів використовує для швидкого, але тимчасового блокування відомого. Це пряме продовження принципу семантичної стійкості (розділи~\ref{chap:ch06}, \ref{chap:ch08}).

Піраміда болю має й стратегічний вимір: вона описує, як виявлення \emph{впливає на економіку} атаки. Блокування гешу змушує зловмисника перекомпілювати зразок --- справа хвилин. Блокування інфраструктури (адрес, доменів) коштує дорожче, але теж посильне. А ось виявлення, що спирається на саму техніку, ставить зловмисника перед вибором: або відмовитися від перевіреного способу дій і розробити новий (дорого, ризиковано), або щоразу спрацьовувати на виявленні. Саме тому мета зрілого виявлення --- не стільки зупинити конкретний зразок, скільки \emph{підняти вартість} усієї операції настільки, щоб вона стала невигідною або помітною. Це перегукується з економічним поглядом на безпеку: абсолютний захист недосяжний, але можна зробити атаку дорожчою за її очікувану вигоду.

Ефективність набору правил зручно характеризувати очікуваним числом виявлених атак. Якщо технікa реалізується з імовірністю $q_i$ у типовій атаці, а правило на неї має повноту $r_i$, то очікувана частка атак, у яких спрацює хоча б одне правило з набору незалежних технік,

\begin{equation}
P_{\text{виявл}} = 1 - \prod_{i}\bigl(1 - q_i r_i\bigr),
\label{eq:ch11-coverage}
\end{equation}

зростає з розширенням покриття технік. Формула~\eqref{eq:ch11-coverage} формалізує цінність \emph{широти} покриття: виявлення, що спирається на кілька незалежних технік атаки, надійніше за будь-яке одиничне правило, бо зловмисник рідко може уникнути одразу всіх поведінкових слідів.

\section{Безфайлові загрози та living-off-the-land}\label{sec:ch11-fileless}

Сучасні атаки дедалі частіше \term{безфайлові}~\cite{ch11:fileless}: код виконується в пам'яті, не залишаючи файлу на диску, а отже невидимий файловому скануванню (розділ~\ref{chap:ch06}). Спорідена тактика --- \term{living-off-the-land} (LotL): зловживання легітимними, вбудованими в систему інструментами (інтерпретаторами команд і сценаріїв, системними утиліами адміністрування), занесеними в проєкт LOLBAS~\cite{ch11:lolbas}. Привабливість для зловмисника: немає стороннього файлу для виявлення, дії зливаються з легітимною адміністративною активністю, обхід контролю застосунків (бо інструменти дозволені).

Виявлення цих загроз неможливе на рівні файлів і спирається на поведінку (розділи~\ref{chap:ch07}, \ref{chap:ch10}): аномальні ланцюги породження процесів, підозрілі командні рядки, нетипове використання легітимних інструментів (наприклад, системна утиліта, що звертається до інтернету й запускає код), перехоплення коду в пам'яті (AMSI, розділ~\ref{chap:ch10}). Виявлення обфускованих сценаріїв застосовує статистичний аналіз самого коду сценарію~\cite{ch11:powershell}: висока ентропія, аномальна структура, рідкісні конструкції. Ключ --- не наявність файлу, а \emph{невідповідність} поведінки нормальному профілю: легітимний інструмент, застосований нетиповим чином і в нетиповому контексті, є сигналом.

Безфайловість і LotL ілюструють глибший зсув у ландшафті загроз: межа між «шкідливим» і «легітимним» дедалі частіше пролягає не у \emph{природі} коду, а в \emph{контексті} його застосування. Інтерпретатор сценаріїв, утиліта віддаленого керування, засіб адміністрування самі по собі законні й потрібні; зловмисними їх робить те, \emph{хто}, \emph{звідки}, \emph{в якій послідовності} й \emph{з якою метою} їх запускає. Це принципово ускладнює виявлення, бо не існує «поганого файлу», який можна занести до чорного списку, --- є лише підозрілі \emph{патерни використання} звичайних засобів. Тому виявлення зміщується від аналізу об'єктів до аналізу \emph{поведінки суб'єктів} і потребує моделі нормальної діяльності, відхилення від якої й породжує сигнал. Саме тут стають незамінними поведінкова аналітика та кореляція (розділ~\ref{chap:ch10}), а також знання типових дій адміністраторів у конкретному середовищі, без якого неможливо відрізнити легітимну адміністративну активність від замаскованої під неї атаки.

Обфускація сценаріїв --- супутня проблема. Зловмисники приховують справжні команди шарами кодування, конкатенації, динамічного складання рядків (розділ~\ref{chap:ch08}), щоб обійти прості правила над текстом команди. Відповідь двоїста: по-перше, перехоплення коду \emph{після} деобфускації на рівні рушія сценаріїв (AMSI, розділ~\ref{chap:ch10}), коли справжні команди вже розгорнуті; по-друге, статистичне виявлення самого факту обфускації~\cite{ch11:powershell} --- аномально висока ентропія, нетипова довжина лексем, рідкісні конструкції, --- бо сильно обфускований адміністративний сценарій сам по собі підозрілий, незалежно від того, що саме він робить.

\section{Життєвий цикл і метрики виявлення}\label{sec:ch11-lifecycle}

Правило виявлення має \term{життєвий цикл}: гіпотеза (яку загрозу й за якими ознаками виявляти) $\to$ розроблення $\to$ тестування (на зразках загрози й на легітимному трафіку) $\to$ розгортання $\to$ моніторинг якості $\to$ налаштування або виведення з експлуатації. Правила старіють: загрози еволюціонують (дрейф, розділ~\ref{chap:ch09}), середовище змінюється, тому несупроводжуване правило з часом або перестає виявляти (низька повнота), або починає шуміти (низька точність).

\begin{tikzfig}
\begin{tikzpicture}[font=\scriptsize,
  b/.style={draw,rounded corners=2pt,minimum height=5.5mm,text width=17mm,align=center,font=\tiny},
  >={Stealth[length=1.4mm]}]
\node[b] (h) at (0,0) {гіпотеза};
\node[b] (d) at (2.2,0) {розроблення};
\node[b] (t) at (4.4,0) {тестування};
\node[b] (dep) at (6.6,0) {розгортання};
\node[b] (m) at (4.4,-1.3) {моніторинг};
\node[b] (tune) at (2.2,-1.3) {налаштування / вивід};
\draw[->] (h)--(d); \draw[->] (d)--(t); \draw[->] (t)--(dep);
\draw[->] (dep)--(m); \draw[->] (m)--(tune); \draw[->] (tune) to[out=120,in=-120] (h);
\end{tikzpicture}
\figcap{Рис.~\thechapter.2 --- Життєвий цикл правила виявлення як безперервний інженерний процес}
\end{tikzfig}

Метрики інженерії виявлення поділяють на якість окремих правил (повнота, точність, хибні тривоги) і характеристики програми виявлення загалом: \emph{покриття} (яку частку технік ATT\&CK покрито правилами), \emph{глибину} (скільки незалежних сигналів на техніку), \emph{час виявлення} (складова dwell time~\eqref{eq:ch10-dwell}), \emph{операційне навантаження} (потік тривог на аналітика). Для оцінювання співвідношення корисного сигналу й шуму застосовують показник на кшталт частки істинних спрацювань серед усіх:

\begin{equation}
\mathrm{PPV} = \frac{TP}{TP + FP},
\label{eq:ch11-ppv}
\end{equation}

де $TP$ --- істинні, $FP$ --- хибні спрацювання. Низький PPV означає, що аналітики витрачають час на хибні тривоги; інженерія виявлення прагне максимізувати PPV за прийнятної повноти, що знову повертає до проблеми базової частоти (розділ~\ref{chap:ch03}) й необхідності кореляції кількох ознак.

Практичний наслідок --- потреба у \emph{зворотному зв'язку} з експлуатації до розроблення. Кожна тривога, яку аналітик позначив як хибну або істинну, --- це розмічений приклад, що дозволяє виміряти реальний PPV правила й скоригувати його. Без такого замкнутого контуру інженер виявлення працює наосліп, оцінюючи правила лише на штучних наборах, які не відображають справжнього розподілу подій у середовищі. Тому зрілі програми інтегрують виявлення з робочим процесом центру реагування (розділ~\ref{chap:ch16}): вердикти аналітиків повертаються як метрики якості правил, живлять їх налаштування й навіть навчальні вибірки для МН-компонентів (розділ~\ref{chap:ch09}). Так інженерія виявлення стає не одноразовим написанням правил, а безперервним циклом, у якому експлуатація й розроблення взаємно живлять одне одного, а якість виявлення вимірюється й покращується систематично, а не інтуїтивно.

\section{Покриття технік і баланс}\label{sec:ch11-coverage}

Зрілі програми виявлення оцінюють і візуалізують \term{покриття} матриці ATT\&CK~\cite{mitre-attack}, виявляючи сліпі зони --- техніки, для яких немає виявлення. Проте покриття не самоціль: важливіша якість виявлення критичних технік, ніж формальне «замальовування» матриці. Оцінки зрілості виявлення~\cite{ch11:detection-maturity} та настанови з організації центрів безпеки~\cite{ch11:alerting} наголошують на балансі широти (покриття), глибини (надійність на техніку) та операційної придатності (прийнятний рівень шуму). Українські практики систематизують розроблення правил для SIEM~\cite{ch11:detection-ua} та формування й застосування індикаторів компрометації~\cite{ch11:ioc-ua}.

Фундаментальний баланс інженерії виявлення --- між повнотою й точністю під операційним обмеженням потоку тривог. Надто агресивні правила дають високу повноту, але топлять аналітиків у хибних тривогах (руйнуючи PPV~\eqref{eq:ch11-ppv}); надто консервативні --- точні, але пропускають. Розв'язання --- не одне «ідеальне» правило, а \emph{набір} правил різної специфічності, кореляція кількох сигналів (розділ~\ref{chap:ch10}) та пріоритезація за ризиком, що дозволяє поєднати широке покриття з керованим потоком тривог.

\section{Кейс: розроблення виявлення для техніки}\label{sec:ch11-case}

Розглянемо розроблення виявлення для техніки зловживання легітимним інструментом (захисна задача, за життєвим циклом підрозділу~\ref{sec:ch11-lifecycle}).

\emph{Крок 1. Гіпотеза.} На основі розвідки загроз (розділ~\ref{chap:ch12}) обираємо техніку ATT\&CK (наприклад, виконання коду через системну утиліту, не призначену для цього) і формулюємо, за якими спостережуваними ознаками вона проявляється.

\emph{Крок 2. Розроблення.} Пишемо хостове правило (Sigma) над телеметрією створення процесів (розділ~\ref{chap:ch10}): легітимна утиліта з нетиповими аргументами, нетиповим батьківським процесом чи мережевою активністю. Орієнтуємося на поведінкові ознаки (вершина піраміди болю), а не на геші.

\emph{Крок 3. Тестування.} Перевіряємо повноту на зразках техніки (зокрема варіаціях) і точність на великому обсязі легітимного трафіку, щоб оцінити хибні тривоги. Налаштовуємо умови для прийнятного PPV~\eqref{eq:ch11-ppv}.

\emph{Крок 4. Розгортання й моніторинг.} Розгортаємо правило, стежимо за якістю в експлуатації, збираємо зворотний зв'язок від аналітиків, коригуємо за появи хибних тривог чи обходів.

\emph{Крок 5. Супровід.} Оновлюємо правило в міру еволюції техніки, документуємо відомі обходи, за втрати ефективності виводимо з експлуатації або замінюємо. Оцінюємо внесок правила в загальне покриття~\eqref{eq:ch11-coverage}.

Кейс ілюструє інженерний характер виявлення: не разовий запис, а керований життєвий цикл із гіпотезою, тестуванням, метриками й супроводом, орієнтований на стійкі поведінкові ознаки та баланс повноти, точності й операційного навантаження.

\section*{Висновки до розділу}
\addcontentsline{toc}{section}{Висновки до розділу}

Інженерія виявлення перетворює методи виявлення на конкретні, супроводжувані, вимірювані правила, розглядаючи їх як код із версіями, тестами, метриками й життєвим циклом. Правила працюють на різних рівнях --- файловому (YARA), хостовому (Sigma), мережевому (Suricata), --- і повне покриття вимагає їх поєднання та кореляції. Піраміда болю (табл.~\thechapter.1) впорядковує індикатори за стійкістю: виявлення на вершині (поведінкові TTP) дороге, але стійке, тоді як геші й адреси дешеві, але нестійкі, тому зріла стратегія інвестує в поведінкові правила~\eqref{eq:ch11-coverage}. Безфайлові й living-off-the-land загрози невидимі на рівні файлів і виявляються лише за поведінкою --- аномальними ланцюгами процесів, командними рядками, перехопленням у пам'яті. Правило має життєвий цикл від гіпотези до виведення з експлуатації й оцінюється метриками якості (повнота, точність, PPV~\eqref{eq:ch11-ppv}) та програми (покриття, глибина, час виявлення, навантаження). Фундаментальний баланс --- між повнотою й точністю під обмеженням потоку тривог --- розв'язують набором правил, кореляцією й пріоритезацією, а не пошуком одного ідеального правила. Головний висновок: ефективне виявлення --- це інженерна дисципліна, орієнтована на стійкі поведінкові ознаки, широке й глибоке покриття технік та кероване операційне навантаження.

\selfcheck
\begin{enumerate}
\item Що таке інженерія виявлення? Чому правила розглядають як код?
\item Назвіть рівні правил виявлення (файлові, хостові, мережеві) та приклади мов. Чому потрібне поєднання рівнів?
\item Які операційні характеристики правила додає інженерія виявлення до метрик детектора?
\item Опишіть піраміду болю (табл.~\thechapter.1). Чому виявлення на вершині стійкіше за виявлення за гешами?
\item Поясніть формулу покриття~\eqref{eq:ch11-coverage}. Чому широта покриття технік підвищує надійність виявлення?
\item Що таке безфайлові атаки та living-off-the-land? Чому вони привабливі для зловмисника?
\item Як виявляють безфайлові та LotL-загрози? Чому ключем є невідповідність поведінки нормальному профілю?
\item Опишіть життєвий цикл правила виявлення. Чому несупроводжуване правило деградує?
\item Які метрики характеризують програму виявлення загалом? Поясніть PPV~\eqref{eq:ch11-ppv} та його зв'язок з навантаженням на аналітиків.
\item Що таке покриття матриці ATT\&CK? Чому формальне покриття не є самоціллю?
\item Як розв'язують баланс повноти й точності під обмеженням потоку тривог?
\item Опишіть розроблення виявлення для техніки (кейс підрозділу~\ref{sec:ch11-case}). Чому орієнтуються на поведінкові ознаки?
\item Як інженерія виявлення пов'язана з розвідкою загроз (розділ~\ref{chap:ch12}) і реагуванням (розділ~\ref{chap:ch16})?
\item Дослідницьке: оцініть покриття обраного набору правил за матрицею ATT\&CK, виявіть сліпі зони й запропонуйте пріоритети їх усунення.
\item Дослідницьке: порівняйте стійкість виявлення на рівні гешів та на рівні поведінкових TTP до варіацій одного сімейства ЗПК.
\end{enumerate}
